\documentclass[10pt,letterpaper]{article}
\usepackage[T1]{fontenc}
\usepackage{lmodern,amsmath,amssymb,microtype}
\usepackage[margin=0.9in]{geometry}
\usepackage[hidelinks]{hyperref}
\providecommand{\tightlist}{\setlength{\itemsep}{0pt}\setlength{\parskip}{0pt}}
\DeclareUnicodeCharacter{2212}{\ensuremath{-}}
\DeclareUnicodeCharacter{2264}{\ensuremath{\le}}
\DeclareUnicodeCharacter{2265}{\ensuremath{\ge}}
\title{An exact marginal sign exclusion below 150 for nine separated subset sums}
\author{Aafi Khan Malek\\ \small Independent researcher, Ahmedabad, India\\
\small\href{mailto:aafimalek2032@gmail.com}{aafimalek2032@gmail.com}}
\date{28 November 2026}
\begin{document}
\maketitle
Fix the central full-sign orientation \texttt{-\/-\/-\/-\/-+++-} in the
nine-real-weight floor-148 branch. We prove that its q50 seven-sign
witness cannot be \texttt{+-\/-\/-\/-++}. This is \textbf{one marginal
sign exclusion for one central orientation}, not a rejection of that
orientation or an exact value of the nine-weight optimum. There are
still 122 unresolved central orientations at the most recent root
screen.

\section{Audited dependency chain}\label{audited-dependency-chain}

The earlier q50 pilot checked all 128 oriented first-seven full-support
words; 87 were rejected exactly. The paired screen then checked all
\(41\cdot128=5248\) residual q50/q51 choices and left just the two q51
partners \texttt{-++++-+} and \texttt{++-\/-\/-++} for the fixed q50
word \texttt{+-\/-\/-\/-++}. The next pilot checked all 128 q53 words
for each of those two choices, so 256 in total. Exact structural checks
cleared both partner fibers except for the single triple of words \[
(v,w,z)=(\texttt{+----++},\ \texttt{-++++-+},\ \texttt{-+---++}).
\] The archived exact replays verify those coverages and the underlying
root precedents. It remains to rule out this single triple; the argument
below does not require any q54 assumption.

\section{Normalize the top triple and keep the
labels}\label{normalize-the-top-triple-and-keep-the-labels}

Write \(Y=x_7,\ X=x_8,\ M=x_9,\ d=M-X\), and use the previously proved
final slice \(148\le X<149\), \(149\le M<150\), \(1\le d<41/39\). Let
\(D_q\) be the q-th smallest positive absolute full nine-sign magnitude
(one per complementary pair). The known rank pattern gives
\(D_{50}=v+d,\ D_{51}=X+M-w\), and \(D_{51}-D_{50}=2g_{158}\), where
\(g_{158}\in[1,2)\) is the actual rank-158 to rank-159 gap. The first
six signs in \(v,w\) are opposites and both seventh signs are plus.
Consequently \(Y=(v+w)/2=X-g_{158}\), and the last three weights form a
short triple, with both adjacent gaps below 2.

Use the order-preserving top-triple normalization proved in the global
block theorem. Its equally spaced replacement at fixed block sum
preserves all labeled subset-sum comparisons, hence preserves the
central word and the words \(v,w,z\) at their rank identities. It weakly
decreases the maximum. The preceding nine-variable lower bound prevents
the new maximum being below 149, so the representative is still in the
floor-148 slice. Relabel it so \(Y=X-1,\ M=X+1\).

In these normalized weights, comparing positive first-six subset
differences with top signed differences 1 and 2 gives first-six
separation at least 3: the two open unit windows cover \((0,3)\). The
five middle top centers \(X-2,X-1,X,X+1,X+2\) give the additional
forbidden interval \((X-3,X+3)\). Therefore adjoining the virtual weight
\(X\) to the first six gives seven \textbf{3-separated} weights.

\section{The sole-survivor
contradiction}\label{the-sole-survivor-contradiction}

Let \(g\) be the central rank gap, so \(1\le g<2\). Our fixed central
word and \(Y=X-1,d=1\) give \[
 g=x_6+X-2-\sum_{i=1}^5x_i,\qquad
 x_6=\sum_{i=1}^5x_i-X+g+2.
\] The first-seven signed form \(a=g+d=g+1\) has word
\(\texttt{-----++}\). The q50 word changes only its first sign, so
\(v=a+2x_1\). The q49/q50 paired rank envelopes give \(98\le v<100\):
\(D_{49}=v-1\ge97\) and \(D_{50}=v+1\le100+(M-149)\). Thus
\(47.5<x_1<49\) and \(2x_1+g\ge97\).

Let \(t=x_2-x_1\) and \(s=x_3-x_2\). The word \(z\) flips exactly
positions 1 and 2 of \(v\), so \(t=(z-v)/2\). The consecutive full slots
50 through 53 have adjacent signed gaps at least 2, giving
\(z-v=D_{53}-D_{50}\ge6\); the companion rank envelopes give \(z<106\),
so \(t\in[3,4)\). The first-six separation already gives \(s\ge3\).

In the virtual 3-separated family \((x_1,\ldots,x_6,X)\), replace the
pair \(x_6,X\) by its positive difference. The resulting six weights
\(x_1,\ldots,x_5,X-x_6\) are still 3-separated: after translating by
\(x_6\), each subset uses exactly one of the replaced pair. The exact
real-six lower bound 24 in
\href{https://arxiv.org/abs/2611.17360}{Notik, arXiv:2611.17360}
therefore implies \(x_5\ge72\) or \(x_6\le X-72<77\). The latter is
impossible: \(x_1>47.5\) and the first-five singleton gaps of at least 3
give \(\sum_{i=1}^5x_i>267.5\), so the exact central relation gives
\(x_6>121.5\). Thus \(x_5\ge72\).

A positive difference \(b\) of two subsets disjoint from \(\{1,2\}\) in
the virtual family obeys \(b\ge3\) and \(|b-t|\ge3\); the latter follows
by comparing opposite translations by \(x_1,x_2\). Because \(t<4\), this
forces \(b\ge t+3\). In particular \(x_4-x_3,X-x_6\ge t+3\). If
\(s\ge6\), summing the resulting first-five lower bounds gives
\(\sum_{i=1}^5x_i\ge4x_1+4t+87\). Together with the central equality and
\(X-x_6\ge t+3\) this implies \[
 4x_1+5t+g+92\le2X.
\] But \(4x_1+g\ge194-g>192\) and \(5t\ge15\), so the left side is above
299 while \(2X<298\). Therefore \(3\le s<6\).

Now consider the virtual block \(x_1,x_2,x_3\), which has short gaps
\(t,s\) at scale 3. Differences of untouched coordinates
\(x_4,x_5,x_6,X\) avoid the radius-3 windows around \(t,s,t+s\), by
opposite translations within this block. The windows overlap, contain 3,
and extend to \(t+s+3\). Thus \(X-x_6\ge t+s+3\). Also
\(x_4\ge x_1+2t+s+3\). Summing with \(x_5\ge72\) and using the central
equation gives \[
 4x_1+5t+3s+80+g\le2X.
\] The distinct first-six subsets \(\{1,6\}\) and \(\{3,4,5\}\) now have
real difference \[
 D=x_1+x_6-(x_3+x_4+x_5)=3x_1+t-X+g+2.
\] The existing bounds give \(D>-1/2\) since \(x_1>47.5\), \(t\ge3\),
\(X<149\) and \(g\ge1\). The preceding displayed inequality gives \[
D\le x_1-\tfrac32 t-\tfrac32 s-38+\tfrac12 g<3,
\] because \(x_1<49\), \(t,s\ge3\), and \(g<2\). This contradicts
first-six separation at least 3 and rules out the last triple of words.

No complete central orientation is removed. On this one central
orientation the remaining q50 words decrease from 40 after the paired
pilot to 39 after this marginal exclusion.

\section{Replay and optional
diagnostics}\label{replay-and-optional-diagnostics}

\begin{verbatim}
sha256sum -c SHA256SUMS
python3 check-complete-marginal.py
\end{verbatim}

The checker reruns the preceding exact coverage and checks the displayed
signed linear arithmetic as exact nonnegative combinations of rational
rows. The written geometric and smaller-dimension premises are not
formalized in Lean. The two optional solver invocations in diagnostics/
returned only numerical survivors of weaker relaxations, including one
other rank slot. They are preserved transparently but are not used for
the proof or counted as exclusions. No additional numerical run was
needed to certify this argument.
\end{document}
